\chapter{Registers}
\label{ch:registers}

Part IV turns from transformation to memory: not the editor's undo history, already covered in Chapter~\ref{ch:undo}, but the several distinct mechanisms Vim provides for holding something so it does not have to be reconstructed from scratch a moment later. Three such mechanisms recur throughout the rest of this book, and it is worth naming the pattern once, here, before treating each separately: a register externalizes a fragment of \textit{text}; a mark, taken up in Chapter~\ref{ch:marks}, externalizes a \textit{position}; and a macro, taken up in Part VI, externalizes a \textit{trajectory of actions}. Each solves the same underlying problem, that working memory is limited and unreliable, at a different grain. This chapter takes the finest grain of the three.

\section{The Unnamed Register}

Every deletion, change, and yank writes its text somewhere, and in the absence of an explicit register name, that somewhere is the unnamed register, addressed as \texttt{"} and read by an unadorned \texttt{p} or \texttt{P}. This is the register a new Vim user interacts with constantly without necessarily knowing its name: \texttt{dw} followed by \texttt{p} moves a word, because the deletion wrote to the unnamed register and the put read from it, with no register name ever typed by the user. The unnamed register is overwritten by nearly every subsequent deleting or yanking operation, which is the single most common source of a specific, frustrating mistake: yanking a piece of text, performing an unrelated deletion in order to navigate somewhere else, and then discovering that the deletion has silently overwritten the text that was meant to be pasted. Named registers, covered next, exist largely to prevent exactly this.

\section{Named Registers}

Twenty-six named registers, addressed \texttt{"a} through \texttt{"z}, hold text independently of the unnamed register and independently of each other. \texttt{"ayy} yanks the current line into register \texttt{a}; \texttt{"ap} puts its contents after the cursor. Because these registers are not touched by ordinary deletions and yanks that do not explicitly name them, they provide exactly the durability the unnamed register lacks: text can be stored in register \texttt{a}, several unrelated edits can be performed, and register \texttt{a}'s contents will still be there when needed.

A register name given in uppercase, rather than lowercase, does not address a different register; it appends to the same one rather than overwriting it. \texttt{"Ayy} appends the current line to whatever register \texttt{a} already contains, rather than replacing it. This is the mechanism behind a common and genuinely useful workflow: collecting several non-adjacent lines from different parts of a buffer into a single register, one appending yank at a time, before putting them all at once somewhere else.

\section{Numbered Registers}

Registers \texttt{"1} through \texttt{"9} are maintained automatically by Vim rather than addressed deliberately by the user in the way named registers are. Register \texttt{"1} holds the most recent deletion of a line or larger span; each subsequent such deletion shifts the previous contents of \texttt{"1} into \texttt{"2}, the previous contents of \texttt{"2} into \texttt{"3}, and so on, up to \texttt{"9}, at which point the oldest entry is discarded. This gives the numbered registers the character of a small, automatic undo history specifically for deleted text, distinct from the undo tree of Chapter~\ref{ch:undo}: \texttt{"3p} retrieves a deletion from three deletions ago, without needing to have named a register at the time of that deletion. Small, sub-line deletions (a single character, or a span smaller than a full line) are handled differently, covered next.

\section{The Small Delete Register}

A deletion smaller than one full line, such as \texttt{x} or \texttt{dw} confined to a single line, is written to the small delete register, addressed \texttt{"-}, rather than to the numbered registers described above. This separation exists so that a long sequence of small, incidental deletions (correcting typos, removing individual characters) does not flood the numbered registers and push out a more significant, larger deletion the user may still want to retrieve. The practical consequence is that \texttt{"1p} is reliable for retrieving a recent line-or-larger deletion specifically, uncontaminated by whatever small corrections happened to be made afterward.

\section{Read-Only Registers}

Several registers are maintained automatically by Vim and cannot be written to directly; they can only be read. \texttt{":} holds the most recently executed Ex command, useful for retrieving and re-executing or adapting a command typed a moment ago without having to retype it from scratch. \texttt{"/} holds the most recent search pattern, which is what makes it possible to insert the current search term directly into a substitution command, covered in Chapter~14, without retyping the pattern. \texttt{".} holds the most recently inserted text, distinct from the repeat command covered in Chapter~\ref{ch:language} in that it holds the literal text itself, retrievable and insertable elsewhere, rather than replaying the insertion as an action. \texttt{"\%} holds the name of the current file, and \texttt{"\#} the name of the alternate file (the buffer's previous file, formalized further in Chapter~\ref{ch:buffers}). These registers are worth knowing exist even if a reader does not use them daily, because each answers a specific, recurring "how do I get Vim to remember this without retyping it" question the moment that question arises.

\section{The Expression Register}

The expression register, addressed \texttt{"=}, is not a storage location in the same sense as the registers above; reading from it evaluates a Vimscript expression, typed at a prompt, and returns the result as text. In Insert mode, \texttt{Ctrl-r =} opens this prompt directly, and typing an expression such as \texttt{5*12} and pressing return inserts the literal text \texttt{60} at the cursor. This is a small preview of the Vimscript material in Part VI, included here because it is, formally, a register like any other: something \texttt{"=p} or \texttt{Ctrl-r =} can read from, even though what it returns is computed rather than stored.

\section{Clipboard Registers}

On a Vim build compiled with clipboard support (confirmed, per Chapter~\ref{ch:installing}, with \texttt{vim --version}), two further registers bridge Vim's internal register system to the operating system's own clipboard. \texttt{"*} corresponds to the primary selection on X11 systems (the text most recently selected with the mouse, pasted with a middle-click) or to the standard clipboard on Windows and macOS; \texttt{"+} corresponds specifically to the standard copy-paste clipboard on systems that distinguish the two (chiefly X11-based Linux desktops). \texttt{"+yy} yanks the current line directly to the system clipboard, from which it can be pasted into any other application; \texttt{"+p} pastes the system clipboard's contents into the buffer.

This is also the appropriate place to note a caution already raised in Chapter~\ref{ch:installing}: clipboard register behavior is the single most environment-dependent piece of the register system, varying with the build of Vim, the terminal emulator, and, on Windows, whether Vim is running under Git Bash, WSL, or natively. A reader whose \texttt{"+y} does not appear to reach the system clipboard should suspect the environment before suspecting the register syntax itself, and confirm clipboard support with \texttt{vim --version} as a first step.

\section{The Black Hole Register}

The black hole register, addressed \texttt{"\_}, discards whatever is written to it and, symmetrically, returns nothing when read. \texttt{"\_dd} deletes the current line without overwriting any other register, unnamed or otherwise, which is the practical use case: performing a deletion purely to remove text, with no intention of ever pasting it back, without that deletion's side effect of clobbering whatever the unnamed register was previously holding. A reader who has been burned once by the unnamed-register overwrite problem described at the start of this chapter will likely find \texttt{"\_d} becoming a reflexive habit for exactly that situation.

\section{Registers as a System, Not a List}

The register catalog in this chapter is long enough to read, at first pass, like an arbitrary collection of special cases: an unnamed default, twenty-six named slots, nine numbered slots with automatic shifting, a small-delete carve-out, several read-only sources, an expression evaluator, two clipboard bridges, and a discard sink. It is worth closing by restating what unifies them: every one of these is addressed with the same \texttt{"}\textit{name} syntax, readable by the same \texttt{p}/\texttt{P} commands and writable, where writable at all, by the same operator-plus-register-prefix syntax covered throughout this chapter. A reader who has learned to yank into register \texttt{a} has, at zero additional syntactic cost, also learned how to yank into the system clipboard, retrieve a three-deletions-ago line, or discard a deletion cleanly, because all four operations are the same grammatical move applied to a different register name. This is the same multiplicative economy Chapter~\ref{ch:language} established for operators and motions, applied here to registers instead.
